\section{Post hoc checks on the replication sets}\label{sec:posthoc}

After the replication we ran four checks on the replication sets, held-out, with the replication's builds, frozen
points, fault rule and instance types. They were planned and frozen with file hashes before they ran
(\texttt{POSTHOCFREEZE.json}), but they come after the replication results and are descriptive.

\paragraph{The six tests with the index preloaded.} The replication's one weak point was the page cache: every fault
rerun was a job of ours, and one pass depended on dropping a faulted block (\S\ref{sec:replication}). We added a product
option that reads the int8 rows and the float16 vectors into anonymous memory at load instead of mapping them
(\texttt{SS\_HNSW\_PRELOAD=1}; result lists byte-identical to the mapped index) and ran the six tests again with the same
plans, seeds and 6 fresh instances per CPU type, every job of ours preloaded. All \nm{pl-n} latency tests pass
(Table~\ref{tab:preload}; ratios \nm{pl-min}--\nm{pl-max}, geometric mean \nm{pl-gm}), with all \nm{pl-blocks} blocks kept,
\nm{pl-reruns} fault reruns and \nm{pl-dropped} dropped blocks. Preloading made our query \nm{pl-slow}\% slower (geometric mean) and
gives up the lazy loading of \S\ref{sec:results}, but the speed result no longer depends on the fault rule. The option stays off by default:
it reads about 2~GB at 1M vectors (int8 rows 0.5~GB, float16 vectors 1.5~GB) before the first query, where the
mapped index loads lazily, and it keeps that memory
anonymous and unreclaimable, where the mapping lets the operating system reclaim pages under memory pressure. The quality
tests read the same lists and give the same results (\nm{pl-qpass} of \nm{pl-n}).

\begin{table}[h]\centering\scriptsize
\caption{Post hoc: the six replication latency tests with our index preloaded into anonymous memory, held-out, same plans
and test as Table~\ref{tab:rep}. Before: the pre-registered ratio of Table~\ref{tab:rep}. Faults: jobs rerun by the fault
rule.}\label{tab:preload}
\setlength{\tabcolsep}{3pt}
\begin{tabular}{lllrrrrrrlr}\toprule
CPU & Set & Rival & Before & Blocks & Ours \us{} & Rival \us{} & Ratio & $p$ & Holm & Faults\\\midrule
\inputtable{tables/rep-preload.tex}
\bottomrule\end{tabular}
\end{table}

\paragraph{Comparison (A) with common kernels and a true int8 bypass.} We ported our SIMD kernels to the rivals' shared C++
path (int8 dot products, float16 widening and 16-row float64 dot products; the retained lists are byte-identical to the
scalar path), added a switch that skips the int8 stage entirely in both harnesses (no int8 reads and no sort; every
candidate goes to the float16 rescoring in generator order), and ran comparison (A) at $C = 400$ on both replication
sets: our walk (preloaded, production adaptive rule) against the rival graphs at their replication builds and FAISS
binary flat with its counting top-$k$, every graph at effort 400 (one instance per CPU type, 5 rounds;
Table~\ref{tab:strong}). The shared stages now cost the rivals \nm{st-shared-min}--\nm{st-shared-max} times what they cost
ours, so the earlier kernel gap is gone; the fastest rival generator that keeps 99\% of exact quality still takes
\nm{st-next-min}--\nm{st-next-max} times as long as ours in total. The int8 screen saves time: without it, rescoring all 400 candidates was slower in \nm{byp-slower} of
\nm{byp-n} system--set--CPU cases (ours \nm{byp-ours-min}--\nm{byp-ours-max}~\us{} slower), with the same retained quality.

\begin{table}[h]\centering\scriptsize
\caption{Post hoc comparison (A), held-out, $C = 400$, common SIMD kernels in the shared stages: \us{} per query (median over
5 rounds; in brackets generator + shared stages) with the int8 screen ($H = 100$) and with no int8 stage ($H = 400$,
``bypass''), for our walk and for the fastest rival generator that keeps 99\% of exact quality; Ratio: rival / ours with
the screen; retained share of exact nDCG@10, ours ; rival.}\label{tab:strong}
\setlength{\tabcolsep}{3pt}
\resizebox{\textwidth}{!}{\begin{tabular}{llrrlrrrl}\toprule
CPU & Set & Ours & Ours, bypass & Rival & Rival & Rival, bypass & Ratio & Retained\\\midrule
\inputtable{tables/strong-a.tex}
\bottomrule\end{tabular}}
\end{table}

\paragraph{Float16 refine for the rivals.} The refining dev-chosen rivals (FAISS HNSW-SQ8 on AMD, FAISS binary HNSW on
Graviton) ran their frozen points with the float16 rows instead of the float32 originals in the refine step
(\nm{f16-n} cases). Their time changed by a factor of \nm{f16-tmin}--\nm{f16-tmax} and their retained share by at most
\nm{f16-qmax}: in these three cases the refine representation does not explain the comparison (B) results.

\paragraph{Where FAISS binary HNSW spends its time.} On MS MARCO dev, with the same codes, $M = 64$ builds and a fixed
$ef = 400$ (our walk without its adaptive rule), \texttt{perf stat} counted each process over the full dev split and over
one query; the difference per query is in Table~\ref{tab:perf}. Our walk computes slightly more Hamming distances than
FAISS (\nm{perf-ndis-ours} against \nm{perf-ndis-faiss} per query), so it does not win by skipping work. FAISS executes
\nm{perf-ins-x-intel}--\nm{perf-ins-x-amd} times as many instructions, has \nm{perf-br-x-intel}--\nm{perf-br-x-arm} times as many branch
misses and takes \nm{perf-cyc-x-intel}--\nm{perf-cyc-x-amd} times as many cycles; about 190 instructions per distance against
50--75 for ours (ours includes its JavaScript call). \texttt{perf record} puts 35--52\% of FAISS's time in its per-candidate
Hamming distance call, 17--23\% in the search loop and about 13--16\% in its heaps. Our walk filters visited nodes without
branches, scores a whole neighbour list in one batch and keeps bucket queues (\S\ref{sec:post}); the counters are consistent with these choices, not with fewer distances, as the source
of the gap; they do not measure the share of each choice. A FAISS setting we did not find could close part of it; the counting top-$k$ of the flat
index (\S\ref{sec:post}) shows that defaults matter, and we will report any correction.

\begin{table}[h]\centering\scriptsize
\caption{Post hoc \texttt{perf stat} per query, MS MARCO 1M dev, same codes, $M = 64$, fixed $ef = 400$: Hamming distances
computed, cycles and instructions (millions), branch misses (thousands), and instructions per distance. Ours includes its
JavaScript call and its shared stages; FAISS ran with the ported SIMD shared path, so its counts are not inflated
by the scalar rescoring.}\label{tab:perf}
\begin{tabular}{llrrrrr}\toprule
CPU & Walk & Distances & Cycles (M) & Instr. (M) & Br. misses (k) & Instr./distance\\\midrule
\inputtable{tables/perf.tex}
\bottomrule\end{tabular}
\end{table}
